% GENERATED FILE. Edit canonical lessons, then run npm run generate:books.
% canonical-source-hash: fnv1a64:b44d8c993d6842d7
% canonical-lessons: TA-C161-vil, TA-C161-koti, TA-C161-varu, TA-C161-arai, TA-C161-kalakku

\chapter{To sell, To boil, To fry, To grind, To mix}
\label{ch:ta-to-sell-to-boil-to-fry-to-grind-to-mix}
\hlchaptermodality{161}

\begin{chapteropening}
\textbf{By the end of this chapter:} \emph{I can say to sell, to boil, to fry, to grind and to mix.}

Complete the last lesson of chapter 161: I can say to sell, to boil, to fry, to grind and to mix.
\end{chapteropening}

\section[\texorpdfstring{\ta{வில்} (vil)}{வில் (vil)}]{\ta{வில்} (vil) — to sell}

\label{lesson:TA-C161-vil}



\begin{quote}
Before the new one: say the Tamil for an advert, then the Tamil for a receipt.
\end{quote}

\subsection*{You'll want to know: \ta{வில்}}
\textbf{\ta{வில்}} — \emph{vil} — \textquotedblleft{}to sell\textquotedblright{}.

\begin{grammarlens}[title={What you've built}]
One more word for this chapter.
\end{grammarlens}

\subsection*{Your turn}
\begin{itemize}
\raggedright
  \item \emph{Say it:} \emph{vil}
  \item \emph{Say it:} \emph{vil}, once more
  \item \emph{Say it:} say \emph{vil}
  \item \emph{Recall:} say the Tamil for an advert, then the Tamil for a receipt, then say \emph{vil} again
\end{itemize}

\begin{tcolorbox}[breakable,colback=teal!4,colframe=teal!35!black,title={Before you move on}]
What does \ta{வில்} mean? (\textquotedblleft{}To sell\textquotedblright{}.) Say it once more.
\end{tcolorbox}
\section[\texorpdfstring{\ta{கொதி} (koti)}{கொதி (koti)}]{\ta{கொதி} (koti) — to boil (of liquids)}

\label{lesson:TA-C161-koti}



\begin{quote}
Before the new one: say the Tamil for a receipt, then the Tamil for to sell.
\end{quote}

\subsection*{You'll want to know: \ta{கொதி}}
\textbf{\ta{கொதி}} — \emph{koti} — \textquotedblleft{}to boil (of liquids)\textquotedblright{}.

\begin{grammarlens}[title={What you've built}]
One more word for this chapter.
\end{grammarlens}

\subsection*{Your turn}
\begin{itemize}
\raggedright
  \item \emph{Say it:} \emph{koti}
  \item \emph{Say it:} \emph{koti}, once more
  \item \emph{Say it:} say \emph{koti}
  \item \emph{Recall:} say the Tamil for a receipt, then the Tamil for to sell, then say \emph{koti} again
\end{itemize}

\begin{tcolorbox}[breakable,colback=teal!4,colframe=teal!35!black,title={Before you move on}]
What does \ta{கொதி} mean? (\textquotedblleft{}To boil (of liquids)\textquotedblright{}.) Say it once more.
\end{tcolorbox}
\section[\texorpdfstring{\ta{வறு} (vaṟu)}{வறு (vaṟu)}]{\ta{வறு} (vaṟu) — to fry, to roast}

\label{lesson:TA-C161-varu}



\begin{quote}
Before the new one: say the Tamil for to sell, then the Tamil for to boil.
\end{quote}

\subsection*{You'll want to know: \ta{வறு}}
\textbf{\ta{வறு}} — \emph{vaṟu} — \textquotedblleft{}to fry, to roast\textquotedblright{}.

\begin{grammarlens}[title={What you've built}]
One more word for this chapter.
\end{grammarlens}

\subsection*{Your turn}
\begin{itemize}
\raggedright
  \item \emph{Say it:} \emph{vaṟu}
  \item \emph{Say it:} \emph{vaṟu}, once more
  \item \emph{Say it:} say \emph{vaṟu}
  \item \emph{Recall:} say the Tamil for to sell, then the Tamil for to boil, then say \emph{vaṟu} again
\end{itemize}

\begin{tcolorbox}[breakable,colback=teal!4,colframe=teal!35!black,title={Before you move on}]
What does \ta{வறு} mean? (\textquotedblleft{}To fry, to roast\textquotedblright{}.) Say it once more.
\end{tcolorbox}
\section[\texorpdfstring{\ta{அரை} (arai)}{அரை (arai)}]{\ta{அரை} (arai) — to grind}

\label{lesson:TA-C161-arai}



\begin{quote}
Before the new one: say the Tamil for to boil, then the Tamil for to fry.
\end{quote}

\subsection*{You'll want to know: \ta{அரை}}
\textbf{\ta{அரை}} — \emph{arai} — \textquotedblleft{}to grind\textquotedblright{}.

\begin{grammarlens}[title={What you've built}]
One more word for this chapter.
\end{grammarlens}

\subsection*{Your turn}
\begin{itemize}
\raggedright
  \item \emph{Say it:} \emph{arai}
  \item \emph{Say it:} \emph{arai}, once more
  \item \emph{Say it:} say \emph{arai}
  \item \emph{Recall:} say the Tamil for to boil, then the Tamil for to fry, then say \emph{arai} again
\end{itemize}

\begin{tcolorbox}[breakable,colback=teal!4,colframe=teal!35!black,title={Before you move on}]
What does \ta{அரை} mean? (\textquotedblleft{}To grind\textquotedblright{}.) Say it once more.
\end{tcolorbox}
\section[\texorpdfstring{\ta{கலக்கு} (kalakku)}{கலக்கு (kalakku)}]{\ta{கலக்கு} (kalakku) — to mix, to stir}

\label{lesson:TA-C161-kalakku}



\begin{quote}
Before the new one: say the Tamil for to fry, then the Tamil for to grind.
\end{quote}

\subsection*{You'll want to know: \ta{கலக்கு}}
\textbf{\ta{கலக்கு}} — \emph{kalakku} — \textquotedblleft{}to mix, to stir\textquotedblright{}.

\begin{grammarlens}[title={What you've built}]
One more word for this chapter.
\end{grammarlens}

\subsection*{Your turn}
\begin{itemize}
\raggedright
  \item \emph{Say it:} \emph{kalakku}
  \item \emph{Say it:} \emph{kalakku}, once more
  \item \emph{Say it:} say \emph{kalakku}
  \item \emph{Recall:} say the Tamil for to fry, then the Tamil for to grind, then say \emph{kalakku} again
\end{itemize}

\begin{tcolorbox}[breakable,colback=teal!4,colframe=teal!35!black,title={Before you move on}]
What does \ta{கலக்கு} mean? (\textquotedblleft{}To mix, to stir\textquotedblright{}.) Say it once more.
\end{tcolorbox}
